%%%%%%%%%%%%%%%%%%%%%%%%%%%%%%%%%%%%%%%%%%%%%%%%%%%%%%%%%%%%%%%%%%%%%%%%%%%%%%
% Conclusion
%
% Summarises the principal findings, research contributions, implications for
% student transition, and directions for future work.
%%%%%%%%%%%%%%%%%%%%%%%%%%%%%%%%%%%%%%%%%%%%%%%%%%%%%%%%%%%%%%%%%%%%%%%%%%%%%%

\section{Conclusion}

This study demonstrated that combining traditional statistical analysis with interpretable machine learning provides a transparent and reproducible framework 
for investigating student transition into first-year university mathematics. Across the evaluated models, measures of prior mathematical preparation 
consistently emerged as the strongest predictors of academic success, with secondary mathematics preparation contributing the greatest overall predictive importance. 
These findings reinforce the importance of discipline-specific educational background in understanding student transition and illustrate how multiple analytical 
approaches can be integrated to provide both predictive performance and educational interpretation.

The study also demonstrates the value of an integrated and reproducible analytical workflow for educational research. 
Statistical modelling, machine learning, independent verification, automated reporting, and publication-ready document generation were combined 
within a single research framework, ensuring that all reported results were derived directly from the underlying analysis. 
This approach reduces opportunities for transcription errors, improves transparency, and provides a reusable foundation for future studies 
investigating student transition, progression, and academic success across comparable educational contexts.

The framework developed in this study provides a foundation for future investigations of student transition using larger cohorts, additional academic disciplines, 
and longitudinal datasets. While the present case study focused on a single first-year mathematics subject, the methodology is readily extensible to other educational 
contexts where transparent interpretation of predictive models is essential. Future work will investigate broader institutional datasets, evaluate additional modelling 
approaches, and further develop the automated research workflow to support reproducible educational analytics and evidence-informed decision making.

\newpage